% C44_POS /root; candidates undergoing independent whole-proof review.
\subsection{C44-C03: characteristic-two conductor strata with a simple pole}
\begin{theorem}[C44-C03]\label{thm:AS-strata}
In characteristic two, fix an Artin--Schreier double-cover pole partition
\[
 d_i=2a_i+1\quad(1\le i\le s),\qquad
 s\ge2,\quad a_i\ge0,\quad a_1=0,
 \qquad g=\sum_i(a_i+1)-1\ge2.
\]
Fix an abelian-fourfold Lefschetz pencil as in H9-T22, and require the
branch locus of the covers to avoid its critical values. At the geometric
generic point of the indicated irreducible conductor stratum, there is
no algebraic class in the new middle trace kernel of the pulled-back
pencil. Consequently
\[
 h_L(z,w)=\tfrac12h_K(\operatorname{cor}z,\operatorname{cor}w),
 \qquad \operatorname{rad}h_L=
 \operatorname{cor}^{-1}(\operatorname{rad}h_K).
 \tag{C44.C03.1}
\]
The identity and vanishing are unconditional. Semipositivity of the
resulting height uses H9-T8 and is conditional on the unrefereed
032/Milne inputs. The assertion includes the geometric generic point
of the hyperelliptic $2$-rank $g-1$ divisor for every $g\ge2$.
It is a statement at geometric generic points, not at every closed
point of the stratum.
\end{theorem}
\begin{proof}
Order the poles and use the dense coordinate chart in which all poles
are finite, with the distinguished simple pole at $t=0$. Put the other
poles at distinct nonzero $b_2,\ldots,b_s$, away from the critical set
$\Sigma$ of the pencil. Choose a rational function $f_0$, regular at zero
and infinity, with exactly these poles and orders $d_2,\ldots,d_s$.
The family
\[
                  w^2-w=f_0(t)+\lambda/t,
                  \qquad \lambda\in\mathbf G_m
 \tag{C44.C03.2}
\]
lies in the specified conductor stratum. Its normalization is smooth
and proper over the parameter. Here is the required local verification
also for a higher pole. If $z$ is a local parameter at a pole of order
$2a+1$, set $v=z^{a+1}w$. The integral local equation becomes
\[
 v^2-z^{a+1}v=z^{2a+2}f(z).
\]
The right side has first term $c z$, with $c$ a unit on the parameter
space. At the unique point $z=v=0$ its derivative in $z$ is $c\ne0$.
The equation is regular there, finite flat of degree two, and agrees
with the extension off the pole, so it is the normalization. Away from
the poles the derivative in $w$ is a unit. This proves the relative
smoothness and justifies smooth proper base change.

Let $V$ be the variable odd local system of the original pencil,
of rank $r>0$, with $q=\#\Sigma$. Its irreducibility and rank-one tame
transvections are the inputs used in H9-T21--T22. Change coordinates
to $x=1/t$, and put
\[
 G=\mathcal L_\psi(f_0(1/x)),\qquad
 F=(\operatorname{inv}^*V)\otimes G.
\]
Their middle extensions to $\mathbf A^1_x$ are irreducible Fourier
sheaves: $G$ has a finite wild pole, and $F$ is irreducible of rank
$r\ge2$. Both are unramified at infinity. The exact character-part
Leray calculation of H9-T22 identifies the new middle cohomology
with
\[
 H^3(A)\otimes\operatorname{NFT}_\psi(G)
       \ \oplus\ \operatorname{NFT}_\psi(F)
 \tag{C44.C03.3}
\]
over the parameter open. This is a calculation of the actual
smooth-proper cohomology, not a substitution of extension by zero at
the old transvections. The outer Leray terms are deck invariant; the
two displayed terms are its nontrivial deck character.

Set $R=\sum_{i=2}^s(d_i+1)$. The Fourier ranks are $R$ and $q+rR$,
by the finite conductor formula, Katz, \emph{Exponential Sums and
Differential Equations}, Lemma 7.3.9(2). Theorem 7.3.8(3) gives
irreducibility of the transforms. At parameter zero, Corollary
7.4.3.1 and Theorem 7.4.4(2)--(3), equivalently Theorem 7.5.4(5)--(7),
give their respective local representations
\[
 J(1,2)\oplus J(1,1)^{R-2},\qquad
 J(1,2)^{r}\oplus J(1,1)^{q+rR-2r}.
 \tag{C44.C03.4}
\]
Each is infinite because it contains a nontrivial unipotent block.
An irreducible representation with infinite image has no nonzero
finite-orbit vector: the span of the orbit would be a nonzero
subrepresentation, and the normal core of its finite-index stabilizer
would then act trivially on the whole representation. This would make
its image finite. Direct sums and tensoring with a constant vector
space preserve the absence of finite-orbit vectors. Thus (C44.C03.3)
has none.

For completeness, the space with ordered pole positions and all
reduced odd principal-part coefficients is irreducible: it is an open
in a vector bundle over the ordered configuration space, with each
leading coefficient nonzero. The above local normalization works
throughout it. It maps dominantly to the indicated conductor stratum
of covers. Its universal trace-kernel local system restricts to
(C44.C03.3). A finite-orbit vector for the universal monodromy would
remain one on this parameter curve. There are therefore no such
vectors at the geometric generic point of the stratum. Algebraic
classes, defined over finite extensions of its function field, have
finite monodromy orbit. The actual trace projector and admissibility
argument of H9-T22 and H9-T24 now give (C44.C03.1).

The genus and $2$-rank are $g$ and $s-1$. Pries--Zhu,
\emph{The $p$-rank stratification of Artin--Schreier curves},
Theorem 1.1 and Corollary 1.3, printed p.~708, identify the components
in characteristic two with partitions of $g+1$ into $s$ positive
parts $a_i+1$, all of dimension $g-1+(s-1)$. For rank $g-1$ the
partition is $(2,1,\ldots,1)$, with $g-1$ parts equal to one.
It has a simple pole and is a divisor in the hyperelliptic moduli
space of dimension $2g-1$. Fixing the target adds three dimensions
to both spaces and leaves codimension unchanged. This proves the
last assertion.
\end{proof}

\subsection{C44-C04: extremal triple-cover splitting strata}
\begin{theorem}[C44-C04]\label{thm:extremal-maroni}
Let $e\ge2$ and $g=3e-2$. In every characteristic, let $U_{e,2e}$
denote the framed space of smooth connected triple covers of the
fixed target $\mathbf P^1$ with Tschirnhausen bundle
$\mathcal O(e)\oplus\mathcal O(2e)$, unramified over the old
critical set. At its geometric generic point the new middle
trace kernel of an abelian-fourfold pencil contains no algebraic
classes. The norm-height identity and radical in (C44.C02.1) hold.
These assertions are unconditional; height semipositivity is
conditional on the unrefereed 032/Milne inputs through H9-T8.

Together with C44-C02 this settles the geometric generic point
of every admissible triple-cover splitting stratum of genus at
least two. It does not assert the same vanishing at every closed
point over $\overline{\mathbf F}_p$.
\end{theorem}
\begin{proof}
The cubic coefficient space for $E=\mathcal O(e)\oplus\mathcal O(2e)$
is $H^0(\mathcal O\oplus\mathcal O(e)\oplus\mathcal O(2e)
\oplus\mathcal O(3e))$. Its smooth connected unramified locus is
irreducible if nonempty. We construct a family in it for each
characteristic and prove that its new trace-kernel monodromy has
no finite-orbit vectors. The restriction argument of C44-C02
will then prove the assertion for this whole stratum.

\paragraph{Characteristic different from three.}
Choose a coordinate $t$ with zero and infinity outside the old critical
set $\Sigma$. Choose $3e-3$ nonzero points $b_i$ so that their
$\mu_3$-orbits are pairwise disjoint and avoid the $\mu_3$-orbits of
$\Sigma$. Put $f(t)=\prod_{i=1}^{3e-3}(t-b_i)$ and consider
\[
                   u^3=(t^3-s)f(t).
 \tag{C44.C04.1}
\]
Remove zero and the finitely many parameters for which a moving root
hits a fixed branch point or a critical value. The normalized covers
are smooth and have $3e$ simple branch points, all with exponent one;
infinity is unramified. Their genus is $3e-2$, and the relative
eigensheaf decomposition is
\[
        \pi_*\mathcal O_C=
        \mathcal O\oplus\mathcal O(-e)\oplus\mathcal O(-2e).
 \tag{C44.C04.2}
\]
The normalized local equations are the tame equations of H9-T23.
Because the degree of the branch polynomial is $3e$, this is a
relative decomposition with fixed line bundles, not only a
fibrewise numerical calculation. It supplies a framing into $U_{e,2e}$.

Fix a nontrivial cubic character $\chi$. On the $t$ line let
$G_\chi=\mathcal L_\chi(f(t))$ and
$F_\chi=V\otimes G_\chi$, always taking middle extensions.
For $\rho:\mathbf A^1_t\to\mathbf A^1_v$, $v=t^3$, set
\[
                  M_\chi=\rho_*G_\chi,
                  N_\chi=\rho_*F_\chi.
\]
Both are irreducible middle extensions. Indeed, restrict to the
$\mu_3$-stable open on which $\rho$ is a Galois cover. Each input
is irreducible. Its three conjugates are pairwise nonisomorphic:
at a fixed branch point one has nontrivial scalar inertia $\chi$,
whereas either conjugate is lisse, by the choice of the $b_i$.
The elementary induction criterion for a normal subgroup with
pairwise inequivalent conjugates therefore makes its induced
representation irreducible. Finite pushforward is its middle
extension, as is checked by inertia invariants at the omitted points.

Finite pushforward and projection formula identify the two new
character terms with
\[
          H^3(A)\otimes\operatorname{MC}_\chi(M_\chi)
                    \ \oplus\ 
                    \operatorname{MC}_\chi(N_\chi).
 \tag{C44.C04.3}
\]
This follows from the actual H9-T23 character cohomology by replacing
$\mathcal L_\chi(t-s)$ there with
$\mathcal L_\chi(t^3-s)$ and then applying $\rho_*$. The fibre
cohomology is the middle convolution, including its local invariant
conditions. Both input sheaves are tame. At infinity their
representations are respectively
\[
                 1\oplus\chi\oplus\chi^2,
                (1\oplus\chi\oplus\chi^2)^{r}.
 \tag{C44.C04.4}
\]
Here the original $V$ and $G_\chi$ are unramified at infinity,
and the ramification of $t\mapsto t^3$ supplies the regular
representation of $\mu_3$.

Katz, \emph{Rigid Local Systems}, Theorem 3.3.3 and Corollary 3.3.6,
give irreducibility of the convolutions and their local invariant
quotient formula. In that formula write
$M_\infty/M_\infty^{I_\infty}$ for the input representation.
Each trivial one-dimensional summand in (C44.C04.4) forces a
$J(1,2)$ block in $M_\infty$. Tensoring by the nontrivial $\chi$
turns it into $J(\chi,2)$, and quotienting by inertia invariants
does not remove it. Thus both outputs have infinite local image.
For a rank check, the first input has finite drops $2$ at zero
and $1$ at each of the $3e-3$ fixed branch images; Corollary
3.3.7 gives output rank $3e-2$. The second has finite drops
$q+(3e-1)r$ in total and output rank $q+(3e-2)r$.
Collisions among old critical images do not change their total
drop; the chosen branch images are disjoint from them.
Irreducibility and infinite image exclude finite-orbit vectors,
as in C44-C03.

\paragraph{Characteristic three.}
Choose $e-1$ distinct nonzero points $b_j$, away from $\Sigma$, so
that the reciprocal points $b_j^{-1}$ and their negatives are
distinct and avoid the reciprocal critical set and its negatives.
For nonzero constants $c_j$, put
\[
 f_0(t)=\sum_{j=1}^{e-1}\frac{c_j}{(t-b_j)^2},
 \qquad w^3-w=f_0(t)+s/t^2,\qquad s\in\mathbf G_m.
 \tag{C44.C04.5}
\]
There are $e$ poles, all of order two, hence genus $3e-2$.
The local order-two normalization and the global basis calculation
of H9-T23--T24, (H9.131) and (H9.144), give a smooth family and
the relative splitting (C44.C04.2). Its pole divisor is fixed,
so once again it maps to the framed stratum.

Put $x=1/t$. For either nonzero additive character index
$\alpha\in\mathbf F_3^\times$, set
\[
 G_\alpha=\mathcal L_\psi(\alpha f_0(1/x)),\qquad
 F_\alpha=(\operatorname{inv}^*V)\otimes G_\alpha,
 \qquad \rho:x\longmapsto v=x^2.
\]
The two finite pushforwards $\rho_*G_\alpha$ and
$\rho_*F_\alpha$ are irreducible by the same induction argument,
now for the tame quadratic cover and its involution $x\mapsto-x$.
A finite wild pole at $b_j^{-1}$ has no pole in its negative orbit,
so the two conjugates cannot be isomorphic. Their ranks are $2$
and $2r$, so neither is an exceptional rank-one additive character.
They are Fourier sheaves. At infinity their representations are
$1\oplus\epsilon$ and $(1\oplus\epsilon)^r$, respectively,
where $\epsilon$ is the nontrivial quadratic tame character.

Projection formula converts the variable term
$\mathcal L_\psi(\alpha s x^2)$ into
$\mathcal L_\psi(\alpha s v)$. The actual new character
cohomology is consequently the sum of $H^3(A)$ tensored with
$\operatorname{NFT}_\psi(\rho_*G_\alpha)$ and
$\operatorname{NFT}_\psi(\rho_*F_\alpha)$, pulled back by
$s\mapsto\alpha s$. Katz, \emph{Exponential Sums and Differential
Equations}, Theorem 7.3.8(3), gives irreducibility. Corollary
7.4.3.1 and Theorem 7.4.4(2)--(3) show that the trivial infinity
summands produce $J(1,2)$ blocks at Fourier parameter zero.
The nontrivial quadratic summands remain nontrivial rank-one
blocks and do not cancel these unipotent blocks. Hence these
transforms also have infinite image and no finite-orbit vectors.

In both characteristics every new character has now been handled.
The direct sum has no finite-orbit vectors. The universal lisse
trace-kernel projector on $U_{e,2e}$ restricts to this sum, so
the restriction argument proves generic algebraic vanishing on
the stratum. The norm-height identity and radical follow by
the same actual-cycle push-pull as in C44-C02. All smooth connected
triple-cover splittings have $e\le f\le2e$; the inequality follows
also directly from the cubic equation, since if $f>2e$ its first
coefficient is zero and the cubic has a linear factor, contradicting
smooth connectedness. C44-C02 handles the strict inequality, while
the present result handles equality for all genus at least two.
\end{proof}

\paragraph{Programme location and residual cover question.}
C44-C03--C04 extend the monodromy method of JAC1964 H9-T21--T24 to
specified special strata. Their height meaning is cote 44, Th.~5,
scan 004, with the primitive qualification of scan 008; the
idempotent $1-b^{-1}g^*g_*$ is the new Chow-motive projector.
The argument concerns lisse sheaves and finite monodromy orbits,
using SGA~7~II, XVIII~6.7.2 for the original pencil variation as
in the predecessor proofs. It gives no positivity theorem for an
extra Tate class at a closed point. Such a class belongs to the
remaining H9-T18 primitive incidence gate, after C44-C01 removes
the positive old surface-axis summand.
